\ncdmodultestmfour{3}{Medier, Tillid og Digital Hverdag}{全真模考卷三：信息批判、政务全数字化与社媒年龄门槛}{社媒新闻完形 / 纸质报刊回归补句 / 儿童图书管理员开放完形 / Digital Post·NemKonto·Cookie·Wi-Fi 等 6 则 / 回信 Jonas 与社媒 16 岁禁令辩论 / 媒体批判与终身学习 Emneopgave 及 Medier og digitalisering 口语讨论}{批判性思维与信源核查 · 严谨因果连接（\textit{som følge heraf, følgelig, hvorved}）· 传闻推测情态（\textit{vist, nok, angiveligt}）· 公共政务数字化词汇}

% =====================================================================
% DEL 1: LÆSNING (50 minutter · 29 point)
% =====================================================================
\ncdsubsec{Del 1: Læsning（阅读理解｜50 minutter · 29 point）}

\begin{tipbox}{考场规则与配速建议（Læsning｜50 minutter · 禁用词典）}
\small
Læsning Opgave 1, 2, 3 og 4 udleveres samtidig. Samlet tid: \textbf{50 minutter}. Kursisten må \textbf{ikke} bruge hjælpemidler（禁用任何词典、笔记或手机）. Der gives 1 point for hvert korrekt svar (maksimum \textbf{29 point}: 8 + 5 + 10 + 6). \textit{Bestået / Ikke bestået} 及格线由各语言中心依 SIRI 指引划定（通常约 60--65\%，建议自测稳定达到 \textbf{20+ point}）。建议配速：Opgave 1（10 min）$\rightarrow$ Opgave 2（12 min）$\rightarrow$ Opgave 3（15 min）$\rightarrow$ Opgave 4（10 min），最后留 3 分钟复查。
\end{tipbox}

% ---------------------------------------------------------------------
\subsubsection*{Opgave 1 (8 point): \textit{Nyheder på sociale medier og kildekritik}}

\begin{emailbox}{Kursistark --- Instruktion（Opgave 1）}
\begin{itemize}
  \item Læs teksten \textit{Nyheder på sociale medier og kildekritik}.
  \item I teksten mangler der otte ord/udtryk (1--8).
  \item I skemaet er der fire ord/udtryk: A, B, C og D.
  \item Sæt kryds i skemaet ved det ord/udtryk, der passer til hvert af hullerne i teksten. Se eksemplet (0).
\end{itemize}

\vspace{0.35em}
\hrule
\vspace{0.45em}

\textbf{\color{NCDNavy}Nyheder på sociale medier og kildekritik}

\vspace{0.3em}
I dag læser stadig flere danskere deres daglige nyheder på sociale medier som Instagram og TikTok i stedet for at købe en avis. Det er nemt og gratis, fordi telefonen altid er i lommen, og nyhederne præsenteres i korte og \textbf{(0)}\ \ldots\ videoer. Men et hurtigt nyhedsfeed skaber også et stort behov for kildekritik.

\vspace{0.25em}
En væsentlig udfordring ved sociale medier er, at algoritmerne først og fremmest viser indhold, som vækker stærke \textbf{(1)}\ \ldots\ som vrede, frygt eller begejstring. Jo flere kommentarer og delinger et opslag får, desto hurtigere spredes det til tusindvis af brugere. \textbf{(2)}\ \ldots\ bliver dramatiske og ensidige påstande ofte vist langt oftere end rolige, nuancerede analyser fra etablerede medier.

\vspace{0.25em}
Derudover er det ikke altid tydeligt, hvem der egentlig står \textbf{(3)}\ \ldots\ et opslag på nettet. Tidligere stod en redaktion med uddannede journalister som garant for, at oplysningerne var \textbf{(4)}\ \ldots, før en artikel blev trykt. I dag kan enhver oprette en flot profil, der ligner et rigtigt nyhedsmedie, og videregive rygter eller manipulerede billeder, som angiveligt stammer fra troværdige kilder.

\vspace{0.25em}
\textbf{(5)}\ \ldots\ er løsningen ikke nødvendigvis at undgå sociale medier helt. Forskere peger på, at man blot skal vænne sig til at bruge få sekunder på at \textbf{(6)}\ \ldots\ sit kritiske blik, før man klikker ''del''. Man kan for eksempel undersøge, om flere uafhængige nyhedsmedier skriver om den samme begivenhed, og se efter, hvem der udtaler sig som ekspert. Hvis en voldsom overskrift mangler dato eller konkrete beviser, bør man være \textbf{(7)}\ \ldots\ indholdet.

\vspace{0.25em}
Flere danske folkeskoler og biblioteker underviser nu børn og unge i digital dannelse og faktatjek, så næste generation bliver rustet til at navigere i strømmen af information. \textbf{(8)}\ \ldots\ handler fremtidens nyhedsvaner ikke om at vælge teknologien fra, men om at stille de rigtige spørgsmål til det, vi ser på skærmen.
\end{emailbox}

\small
\renewcommand{\arraystretch}{1.28}
\begin{longtable}{@{\extracolsep{\fill}} >{\bfseries\color{NCDNavy}\centering\arraybackslash}p{0.8cm} >{\raggedright\arraybackslash}p{3.45cm} >{\raggedright\arraybackslash}p{3.45cm} >{\raggedright\arraybackslash}p{3.45cm} >{\raggedright\arraybackslash}p{3.45cm} @{}}
\toprule
\textbf{Nr.} & \textbf{A} & \textbf{B} & \textbf{C} & \textbf{D} \\
\midrule
\endfirsthead
\toprule
\textbf{Nr.} & \textbf{A} & \textbf{B} & \textbf{C} & \textbf{D} \\
\midrule
\endhead
\bottomrule
\endfoot
\textcolor{NCDGray}{0} & \textcolor{NCDGray}{komplicerede} \hfill \fbox{\phantom{X}} & \textcolor{NCDGray}{dyre} \hfill \fbox{\phantom{X}} & \textcolor{NCDGray}{fængende} \hfill \fbox{\textbf{X}} & \textcolor{NCDGray}{gamle} \hfill \fbox{\phantom{X}} \\
\midrule
1 & regler \hfill \fbox{\phantom{X}} & følelser \hfill \fbox{\phantom{X}} & vaner \hfill \fbox{\phantom{X}} & udgifter \hfill \fbox{\phantom{X}} \\
\midrule
2 & Derfor \hfill \fbox{\phantom{X}} & Tværtimod \hfill \fbox{\phantom{X}} & Til gengæld \hfill \fbox{\phantom{X}} & Trods det \hfill \fbox{\phantom{X}} \\
\midrule
3 & ved \hfill \fbox{\phantom{X}} & mellem \hfill \fbox{\phantom{X}} & bag \hfill \fbox{\phantom{X}} & over \hfill \fbox{\phantom{X}} \\
\midrule
4 & opdigtede \hfill \fbox{\phantom{X}} & ensidige \hfill \fbox{\phantom{X}} & tilfældige \hfill \fbox{\phantom{X}} & dokumenterede \hfill \fbox{\phantom{X}} \\
\midrule
5 & Alligevel \hfill \fbox{\phantom{X}} & Følgelig \hfill \fbox{\phantom{X}} & Desværre \hfill \fbox{\phantom{X}} & Tidligere \hfill \fbox{\phantom{X}} \\
\midrule
6 & glemme \hfill \fbox{\phantom{X}} & skærpe \hfill \fbox{\phantom{X}} & skjule \hfill \fbox{\phantom{X}} & opgive \hfill \fbox{\phantom{X}} \\
\midrule
7 & tilfreds med \hfill \fbox{\phantom{X}} & afhængig af \hfill \fbox{\phantom{X}} & skeptisk over for \hfill \fbox{\phantom{X}} & glad for \hfill \fbox{\phantom{X}} \\
\midrule
8 & Tværtimod \hfill \fbox{\phantom{X}} & For eksempel \hfill \fbox{\phantom{X}} & Desværre \hfill \fbox{\phantom{X}} & I sidste ende \hfill \fbox{\phantom{X}} \\
\end{longtable}
\normalsize

% ---------------------------------------------------------------------
\subsubsection*{Opgave 2 (5 point): \textit{Avisen der gik tilbage til papir}}

\begin{emailbox}{Kursistark --- Instruktion（Opgave 2）}
\begin{itemize}
  \item Læs teksten \textit{Avisen der gik tilbage til papir}.
  \item I hvert afsnit er der et hul, hvor der mangler en sætning.
  \item Sætningerne, der mangler, står på listen (A--H).
  \item Find den sætning, der passer bedst i hvert afsnit (1--5).
  \item Skriv dit svar (= bogstav) på linjen i hullet.
  \item Der er \textbf{to sætninger, du ikke skal bruge}. Se eksemplet i første afsnit.
\end{itemize}

\vspace{0.35em}
\hrule
\vspace{0.45em}

\textbf{\color{NCDNavy}Avisen der gik tilbage til papir}

\vspace{0.3em}
Da lokalavisen \textit{Fjordposten} i en mindre jysk kystby for fem år siden besluttede at lukke sin trykte avis og udkomme rent digitalt, lød det som den eneste fornuftige vej frem. Trykkeriet og distributionen var blevet alt for dyre i forhold til annonceindtægterne. \textbf{(0)}\ \underline{\hspace{0.45cm}\textbf{A}\hspace{0.45cm}}\textbf{.} Men allerede efter to år begyndte redaktør Mette Krogh og hendes journalister at mærke en uventet udvikling, som fik dem til at tvivle på beslutningen.

\vspace{0.25em}
Selvom hjemmesiden stadig havde en del hurtige klik, faldt den samlede læsetid dramatisk. Mange lokale læsere åbnede kun en artikel i ti sekunder via Facebook, før de scrollede videre til noget andet. \textbf{(1)}\ \rule{1.8cm}{0.4pt}. Samtidig skrev flere ældre abonnenter til redaktionen, at de savnede at sidde med en fysisk avis ved morgenbordet lørdag formiddag uden skærme og blinkende reklamer.

\vspace{0.25em}
Efter længere overvejelser tog redaktionen et skridt, som mange i mediebranchen kaldte gammeldags. De valgte at relancere \textit{Fjordposten} som en tyk avis på kraftigt papir hver anden uge, fyldt med længere reportager og portrætter fra lokalområdet. \textbf{(2)}\ \rule{1.8cm}{0.4pt}. I stedet for at jage hurtige nyheder om trafikulykker ville avisen nu fokusere på baggrundsartikler, debatindlæg og fortællinger om lokale ildsjæle, som man ikke finder andre steder.

\vspace{0.25em}
Alligevel var starten økonomisk sårbar, for trykkepriser og papirudgifter var steget kraftigt i de senere år. Redaktionen måtte tænke kreativt for at få budgettet til at hænge sammen uden at gå på kompromis med kvaliteten. \textbf{(3)}\ \rule{1.8cm}{0.4pt}. Det viste sig heldigvis at være en succes, fordi byens bagere, boghandlere og kulturhuse gerne ville støtte et lokalt medie, hvor deres annoncer faktisk blev læst i ro og mag.

\vspace{0.25em}
I dag har \textit{Fjordposten} fået flere hundrede nye abonnenter, herunder overraskende mange børnefamilier i trediverne. Flere fortæller til redaktionen, at avisen bliver liggende på spisebordet hele weekenden, så både voksne og større børn bladrer i den. \textbf{(4)}\ \rule{1.8cm}{0.4pt}. På den måde skaber papiret samtaler ved middagsbordet om byens nye svømmehal, skoler og foreninger, hvilket sjældent sker, når alle sidder med hver sin mobilskærm.

\vspace{0.25em}
Redaktør Mette Krogh tror dog ikke, at alle danske aviser bør slukke deres digitale platforme. Papir og internet udelukker nemlig ikke hinanden, men kan tværtimod styrke den lokale journalistik sammen. \textbf{(5)}\ \rule{1.8cm}{0.4pt}. ''For os har papiret ikke været et skridt bagud, men en påmindelse om, at god journalistik kræver tid --- både for dem, der skriver den, og for dem, der læser den,'' afslutter hun.
\end{emailbox}

\begin{grammatikbox}{Liste (A--H)（A 为例句 · 5 项作答 + 2 项多余干扰项）}
\small
\begin{itemize}[leftmargin=2.2em, itemsep=0.35em]
  \item[\textbf{\color{NCDGray}A.}] \textcolor{NCDGray}{Samtidig læste stadig flere yngre borgere nyheder på deres smartphone i bussen eller i frokostpausen.} {\footnotesize\color{NCDRed}[Eksempel (0)]}
  \item[\textbf{B.}] Derudover læser de fleste danskere nu udelukkende engelske aviser om internationale nyheder.
  \item[\textbf{C.}] Planen var aldrig at konkurrere med internettet på hurtighed, men derimod at skabe et pusterum til fordybelse.
  \item[\textbf{D.}] Mens nettet leverer hurtige opdateringer i hverdagen, giver den trykte udgave tid til eftertanke og tillid i lokalsamfundet.
  \item[\textbf{E.}] I modsætning til en lukket telefon kan en avis deles fysisk mellem familiens medlemmer over flere dage.
  \item[\textbf{F.}] Derfor faldt antallet af betalende abonnenter, fordi folk ikke ville betale for korte tekster, de hurtigt glemte igen.
  \item[\textbf{G.}] Alligevel besluttede kommunen at lukke byens folkebibliotek for at spare penge på elregningen.
  \item[\textbf{H.}] Derfor indgik de et tæt samarbejde med lokale butikker om både annoncer og salg af avisen ved kassen.
\end{itemize}
\end{grammatikbox}

% ---------------------------------------------------------------------
\subsubsection*{Opgave 3 (10 point): \textit{Nadia arbejder som børnebibliotekar}（无词库开放完形）}

\begin{emailbox}{Kursistark --- Instruktion（Opgave 3）}
\begin{itemize}
  \item Læs teksten \textit{Nadia arbejder som børnebibliotekar}. Der mangler 10 ord (1--10).
  \item Skriv et ord, der passer. Du skal kun skrive \textbf{ét ord}.
\end{itemize}

\vspace{0.35em}
\hrule
\vspace{0.45em}

\textbf{\color{NCDNavy}Nadia arbejder som børnebibliotekar}

\vspace{0.3em}
Nadia har arbejdet som børnebibliotekar på et stort folkebibliotek i Aarhus i syv år. Hendes vigtigste opgave er at vække børns glæde ved historier i en tid, \textbf{(1)}\ \rule{1.8cm}{0.4pt}\ skærme fylder mere end nogensinde. Selvom mange voksne tror, at børn slet \textbf{(2)}\ \rule{1.8cm}{0.4pt}\ gider åbne en fysisk bog længere, oplever Nadia hver uge lange køer ved udlånsskranken.

\vspace{0.25em}
Om formiddagen får biblioteket ofte besøg af børnehaver og skoleklasser. Nadia sidder på puder på gulvet sammen \textbf{(3)}\ \rule{1.8cm}{0.4pt}\ børnene og læser højt fra eventyr og helt nye billedbøger. Bagefter snakker de om billederne og stiller spørgsmål. For Nadia handler arbejdet nemlig lige så meget om nærvær \textbf{(4)}\ \rule{1.8cm}{0.4pt}\ om selve bøgerne.

\vspace{0.25em}
Ud over højtlæsningen hjælper Nadia forældre med at finde gode bøger til deres børns alder. Nogle børn har svært \textbf{(5)}\ \rule{1.8cm}{0.4pt}\ at komme i gang med læsningen, og så finder hun sjove tegneserier om emner, \textbf{(6)}\ \rule{1.8cm}{0.4pt}\ interesserer barnet. \textbf{(7)}\ \rule{1.8cm}{0.4pt}\ arrangerer hun små bogklubber og værksteder i weekenderne, så familierne mødes gratis.

\vspace{0.25em}
Det sværeste ved jobbet er, at bibliotekets budget flere gange er blevet skåret ned. For eksempel mangler der ofte penge \textbf{(8)}\ \rule{1.8cm}{0.4pt}\ at købe nok nye børnebøger. Alligevel er kollegaerne gode til at støtte hinanden, når der skal planlægges aktiviteter, \textbf{(9)}\ \rule{1.8cm}{0.4pt}\ de travle perioder før sommerferien bliver nemmere at overskue.

\vspace{0.25em}
Når Nadia ser et barn gå glad hjem med en hel stak bøger under armen, ved hun præcis, hvorfor hun elsker sit arbejde. Hun er overbevist \textbf{(10)}\ \rule{1.8cm}{0.4pt}, at et stærkt børnebibliotek giver alle børn lige muligheder --- helt uafhængigt af forældrenes pengepung.
\end{emailbox}

% ---------------------------------------------------------------------
\subsubsection*{Opgave 4 (6 point): 6 korte tekster（数字政务、隐私安全与校园日常）}

\begin{emailbox}{Kursistark --- Instruktion（Opgave 4）}
\begin{itemize}
  \item Læs de seks små tekster (1--6).
  \item Sæt \textbf{ét kryds} for det rigtige svar (A--C).
\end{itemize}
\end{emailbox}

\begin{emailbox}{1. Digital Post og fritagelse}
\textbf{Digital Post og fritagelse}\\
Alle borgere i Danmark over 15 år modtager som udgangspunkt breve fra offentlige myndigheder i Digital Post, som du kan læse med MitID via borger.dk, e-Boks eller mit.dk. Hvis du på grund af sygdom, funktionsnedsættelse, sprogvanskeligheder eller manglende adgang til internet i hjemmet ikke kan bruge digital selvbetjening, kan du søge om at blive fritaget. Ansøgningen sker ved personligt fremmøde i kommunens Borgerservice, eller ved at en pårørende møder op med en skriftlig fuldmagt. Bliver du fritaget, får du fremover offentlige breve sendt med fysisk post.
\hfill {\small\textit{Kilde: borger.dk, 2026 (red.)}}

\vspace{0.3em}
\hrule
\vspace{0.3em}
\small
\fbox{\phantom{X}}\ \textbf{A.} Alle borgere bliver automatisk fritaget fra Digital Post, når de fylder 65 år.\\[0.25em]
\fbox{\phantom{X}}\ \textbf{B.} Borgere, der bliver fritaget fra Digital Post, modtager i stedet vigtige breve med posten.\\[0.25em]
\fbox{\phantom{X}}\ \textbf{C.} Man må kun søge om fritagelse fra Digital Post digitalt via e-Boks.
\end{emailbox}

\begin{emailbox}{2. Ændring af NemKonto i banken}
\textbf{Ændring af NemKonto i banken}\\
Alle borgere over 18 år i Danmark skal have en NemKonto, hvor det offentlige kan udbetale for eksempel overskydende skat, SU, løn eller boligstøtte. Hvis du skifter bank eller ønsker at bruge en anden konto som din NemKonto, kan du selv ændre den på nemkonto.dk med MitID eller bede din nye bank om at registrere kontoen for dig. Bemærk, at der kan gå et par hverdage, før ændringen træder i kraft. Derfor bør du altid vente med at lukke din gamle bankkonto, indtil din nye NemKonto er aktiv, så udbetalinger ikke bliver forsinkede.
\hfill {\small\textit{Kilde: nemkonto.dk, 2026 (red.)}}

\vspace{0.3em}
\hrule
\vspace{0.3em}
\small
\fbox{\phantom{X}}\ \textbf{A.} Offentlige myndigheder bruger kun NemKonto til at udbetale skat, ikke SU eller løn.\\[0.25em]
\fbox{\phantom{X}}\ \textbf{B.} Man kan ikke ændre sin NemKonto selv, men skal altid kontakte kommunens Borgerservice.\\[0.25em]
\fbox{\phantom{X}}\ \textbf{C.} Du skal vente med at lukke din gamle konto, til din nye NemKonto er trådt i kraft.
\end{emailbox}

\begin{emailbox}{3. Kort SMS fra kulturhuset om flyttet aftenforedrag}
Hej Astrid\\
Desværre er vores oplægsholder til torsdagens aftenforedrag om falske nyheder blevet syg. Foredraget er derfor flyttet til næste tirsdag kl. 19.00 i samme sal. Din billet gælder automatisk til den nye dato, men hvis du ikke kan komme der, kan du få hele beløbet refunderet ved at svare på denne besked inden mandag kl. 12.\\
Dbh Kulturhuset Vesterbro

\vspace{0.3em}
\hrule
\vspace{0.3em}
\small
\fbox{\phantom{X}}\ \textbf{A.} Aftenforedraget flyttes til næste uge, og man kan få billetten refunderet, hvis datoen ikke passer.\\[0.25em]
\fbox{\phantom{X}}\ \textbf{B.} Alle gæster skal købe en helt ny billet for at komme ind til foredraget næste tirsdag.\\[0.25em]
\fbox{\phantom{X}}\ \textbf{C.} Man kan først bede om at få pengene tilbage efter foredraget på tirsdag.
\end{emailbox}

\begin{emailbox}{4. Cookies og databeskyttelse på nettet}
\textbf{Cookies og databeskyttelse på nettet}\\
Når du besøger en dansk hjemmeside, må den ikke indsamle oplysninger om din adfærd til statistik eller målrettet markedsføring uden dit klare samtykke. Hjemmesiden skal give dig et reelt valg, så det er lige så nemt at afvise disse cookies som at acceptere dem. Kun teknisk nødvendige cookies, som får indkøbskurven eller login til at fungere, må bruges uden samtykke. Hvis du siger nej til statistik- og markedsføringscookies, har du stadig ret til at læse hjemmesidens almindelige indhold, og du kan til enhver tid ændre dit valg igen.
\hfill {\small\textit{Kilde: datatilsynet.dk, 2026 (red.)}}

\vspace{0.3em}
\hrule
\vspace{0.3em}
\small
\fbox{\phantom{X}}\ \textbf{A.} Hjemmesider skal altid bede om dit samtykke, før de bruger teknisk nødvendige cookies til login.\\[0.25em]
\fbox{\phantom{X}}\ \textbf{B.} Du har ret til at afvise markedsføringscookies og alligevel bruge hjemmesidens almindelige indhold.\\[0.25em]
\fbox{\phantom{X}}\ \textbf{C.} Har du først sagt ja til alle cookies én gang, kan du ikke trække dit samtykke tilbage senere.
\end{emailbox}

\begin{emailbox}{5. Kort SMS fra studiekammerat om lærebog}
Hej Tobias\\
Har du stadig brug for vores lærebog i samfundsfag til opgaven næste uge? Jeg er selv færdig med mit projekt nu, så du må meget gerne låne bogen fra i morgen eftermiddag, hvis du ikke allerede har købt den. Jeg kan tage den med i kantinen kl. 14, hvis du lige skriver tilbage i aften.\\
Kh Emilie

\vspace{0.3em}
\hrule
\vspace{0.3em}
\small
\fbox{\phantom{X}}\ \textbf{A.} Emilie vil gerne låne Tobias' lærebog til sit eget projekt i morgen.\\[0.25em]
\fbox{\phantom{X}}\ \textbf{B.} Emilie spørger, om Tobias kan købe en ny lærebog til hende i kantinen.\\[0.25em]
\fbox{\phantom{X}}\ \textbf{C.} Emilie tilbyder at låne sin lærebog ud til Tobias fra i morgen eftermiddag.
\end{emailbox}

\begin{emailbox}{6. Sikkerhed ved offentligt wi-fi og MitID}
\textbf{Sikkerhed ved offentligt wi-fi og MitID}\\
Når du sidder på en café, i toget eller i lufthavnen, kan et åbent, offentligt wi-fi-netværk være praktisk, men forbindelsen er ofte ikke krypteret. Derfor frarådes det, at du logger på din netbank, tjekker følsomme sundhedsdata eller godkender betalinger med MitID over et ukendt åbent netværk. Brug i stedet din telefons eget mobildata eller en sikker VPN-forbindelse, hvis du skal ordne vigtige digitale ærinder på farten. Husk også altid at tjekke teksten i MitID-appen grundigt, før du swiper, og afvis straks anmodninger, du ikke selv har startet.
\hfill {\small\textit{Kilde: sikkerdigital.dk, 2026 (red.)}}

\vspace{0.3em}
\hrule
\vspace{0.3em}
\small
\fbox{\phantom{X}}\ \textbf{A.} Du bør vælge mobildata i stedet for åbent wi-fi, når du skal ordne bankærinder med MitID.\\[0.25em]
\fbox{\phantom{X}}\ \textbf{B.} Offentlige wi-fi-netværk i tog og på caféer er altid krypterede og lige så sikre som mobildata.\\[0.25em]
\fbox{\phantom{X}}\ \textbf{C.} Man skal altid bruge MitID-appen, hver gang man vil koble sin telefon til et nyt wi-fi-netværk.
\end{emailbox}

% =====================================================================
% DEL 2: SKRIVNING (Dual Tasks)
% =====================================================================
\ncdsubsec{Del 2: Skrivning（书面表达｜Opgave 1 私人回信 \& Opgave 2 论坛辩论贴）}

\subsubsection*{Opgave 1（私人邮件回复 · Minimum 125 ord）}

\textbf{Situation:} Du har fået en e-mail fra din ven, Jonas.\\
\textbf{Opgave:} Læs e-mailen, og skriv et svar til Jonas. Du skal svare på Jonas' spørgsmål.\\
\textbf{Længde:} Minimum 125 ord.

\begin{emailbox}{Modtaget e-mail fra Jonas}
Hej.

\vspace{0.3em}
Hvordan går det? Det er efterhånden længe siden, vi sidst snakkede sammen. Sidst vi mødtes, fortalte du jo, at du havde slået alle nyhedsnotifikationer fra på telefonen om aftenen, fordi du blev stresset af alle de dårlige nyheder. Hvordan er det gået med det? Føler du, at du sover bedre, eller er du bange for at gå glip af noget vigtigt?

\vspace{0.3em}
Her hos mig går det fint, men jeg vil også gerne have bedre nyhedsvaner. Jeg overvejer enten at abonnere på en trykt weekendavis, som jeg kan læse i ro og mag, eller at lytte til gratis daglige podcasts om nyheder på vej til arbejde. Jeg er lidt i tvivl om, hvad jeg skal vælge. Hvad synes du?

\vspace{0.3em}
Forresten så jeg i går, at min søster lægger rigtig mange billeder af sin lille søn ud på sociale medier. Det fik mig til at tænke over børns privatliv på nettet. Synes du, at forældre bør dele billeder af små børn på sociale medier, eller bør de lade være? Hvorfor?

\vspace{0.3em}
\ldots
\end{emailbox}

\begin{emailbox}{Dit svar｜Til: Jonas \qquad Emne: Sv: Nyhedsvaner, avis eller podcast og børns privatliv}
\vspace{0.3em}
\rule{\linewidth}{0.4pt}\\[0.55em]
\rule{\linewidth}{0.4pt}\\[0.55em]
\rule{\linewidth}{0.4pt}\\[0.55em]
\rule{\linewidth}{0.4pt}\\[0.55em]
\rule{\linewidth}{0.4pt}\\[0.55em]
\rule{\linewidth}{0.4pt}\\[0.55em]
\rule{\linewidth}{0.4pt}\\[0.55em]
\rule{\linewidth}{0.4pt}\\[0.55em]
\rule{\linewidth}{0.4pt}\\[0.55em]
\rule{\linewidth}{0.4pt}\\[0.55em]
\rule{\linewidth}{0.4pt}\\[0.55em]
\rule{\linewidth}{0.4pt}
\end{emailbox}

\subsubsection*{Opgave 2（观点辩论帖 / \textit{Blogindlæg} · Minimum 100 ord）}

\textbf{Situation:} Søren Holm har skrevet et blogindlæg på en hjemmeside, hvor man diskuterer aktuelle emner.\\
\textbf{Opgave:} Skriv et blogindlæg, hvor du diskuterer Søren Holms synspunkt om, at der bør indføres en aldersgrænse på 16 år for sociale medier. Er du enig? Hvorfor? / Hvorfor ikke?\\
\textbf{Længde:} Minimum 100 ord.

\begin{emailbox}{Blogindlæg af Søren Holm：\textit{Indfør en aldersgrænse på 16 år for sociale medier}}
\textbf{\color{NCDNavy}Indfør en aldersgrænse på 16 år for sociale medier}

\vspace{0.3em}
Det er på høje tid, at vi beskytter børn og unge bedre mod sociale medier som TikTok og Instagram. I dag bruger mange 11- til 15-årige flere timer hver aften på at scrolle gennem korte videoer, hvilket fører til dårlig søvn, koncentrationsbesvær i skolen og et usundt pres for at se perfekt ud. Selvom platformene officielt har en aldersgrænse på 13 år, er den alt for nem at omgå, og digital mobning spredes lynhurtigt blandt unge. Derfor bør politikerne indføre en klar og streng aldersgrænse på 16 år ved lov, så børn først får adgang, når de er mere modne.
\end{emailbox}

\begin{emailbox}{Dit blogindlæg｜Diskussion af Søren Holms synspunkt}
\vspace{0.3em}
\rule{\linewidth}{0.4pt}\\[0.55em]
\rule{\linewidth}{0.4pt}\\[0.55em]
\rule{\linewidth}{0.4pt}\\[0.55em]
\rule{\linewidth}{0.4pt}\\[0.55em]
\rule{\linewidth}{0.4pt}\\[0.55em]
\rule{\linewidth}{0.4pt}\\[0.55em]
\rule{\linewidth}{0.4pt}\\[0.55em]
\rule{\linewidth}{0.4pt}\\[0.55em]
\rule{\linewidth}{0.4pt}\\[0.55em]
\rule{\linewidth}{0.4pt}\\[0.55em]
\rule{\linewidth}{0.4pt}
\end{emailbox}

% =====================================================================
% DEL 3: MUNDTLIG KOMMUNIKATION
% =====================================================================
\ncdsubsec{Del 3: Mundtlig kommunikation（口语交际全真题卡｜Opgave 1 \& 2）}

\subsubsection*{Mundtlig Opgave 1 (Emneopgave｜Del 1: 1--2 min. + Del 2: 3--4 min.)}

\begin{tipbox}{DU3 Modul 4 官方口试规则（SIRI / Modultest.dk \textit{Opgave 1: Emneopgave}）}
\small
\begin{itemize}
  \item \textbf{Del 1}：\textit{Du vælger ved lodtrækning ét af de to emner, som du har forberedt. Herefter fortæller du om emnet i 1--2 minutter.}（考生须事先围绕 \textbf{danske forhold} 准备两个主题，考场抽签决定其一并脱稿连贯陈述 1--2 分钟。）
  \item \textbf{Del 2}：\textit{Præsentationen følges op af spørgsmål til emnet i 3--4 minutter. Tester stiller spørgsmål, der lægger op til, at du forklarer, uddyber og begrunder i lidt længere taleture, fx: Vil du fortælle lidt mere om …? / Kan du give nogen eksempler på …? / Hvad er fordelene / ulemperne ved …? Hvorfor? / Hvad synes du om …? Hvorfor?}
  \item \textbf{纯关键词（stikord）铁律}：官方明确规定只能携带只含单词/短语的提示卡（\textit{stikord, dvs. enkeltord og ikke hele sætninger}），严禁写完整句子照读；且内容必须同时兼顾\textbf{宏观丹麦社会现状（generelt fokus）}与\textbf{个人亲身体验（egne erfaringer）}，整体评估为 \textit{Bestået / Ikke bestået}。
\end{itemize}
\end{tipbox}

\begin{grammatikbox}{Stikord-kort 3A（已填纯关键词卡）：Moderne livsstil: nyhedsvaner og kildekritik}
\small
\renewcommand{\arraystretch}{1.3}
\begin{tabularx}{\linewidth}{@{} >{\leavevmode\bfseries\color{NCDNavy}\raggedright\arraybackslash}p{4.3cm} >{\raggedright\arraybackslash}X @{}}
\toprule
\multicolumn{2}{c}{\normalsize\bfseries\color{NCDRed} Emne 3A: Moderne livsstil med fokus på nyhedsvaner og kildekritik \quad\small\color{NCDGray}（仅含 stikord）} \\
\midrule
1. Generelt fokus（丹麦新闻生态） & DR og TV 2 (public service) · sociale medier · push-beskeder · algoritmer \\
\midrule
2. Fordele ved digitale medier & hurtigt overblik · gratis nyheder · mange perspektiver · nem adgang på farten \\
\midrule
3. Ulemper / kildekritik & falske nyheder · clickbait · ekkokamre · manglende kildetjek · informationsstress \\
\midrule
4. Egne erfaringer \& vaner & sammenligne to uafhængige kilder · slå notifikationer fra · rolig weekendavis \\
\midrule
5. Vurdering（个人立场） & digital dannelse i skolen · bevidst pause fra telefonen om aftenen \\
\bottomrule
\end{tabularx}

\vspace{0.45em}
\textbf{\color{NCDRed}Del 2 考官追问预测（Eksaminators opfølgende spørgsmål · 4 题）：}
\begin{enumerate}
  \item \textit{Vil du fortælle lidt mere om, hvordan du selv vurderer, om en nyhed på internettet er troværdig?}
  \item \textit{Kan du give nogle eksempler på, hvordan nyhedsstrømmen på sociale medier påvirker din egen hverdag?}
  \item \textit{Hvad er efter din mening fordelene og ulemperne ved, at de fleste læser nyheder gratis på mobilen i stedet for i en trykt avis?}
  \item \textit{Hvad synes du om den rolle, som public service-medier som DR og TV 2 spiller i Danmark? Hvorfor?}
\end{enumerate}
\end{grammatikbox}

\begin{grammatikbox}{Stikord-kort 3B（已填纯关键词卡）：Uddannelse: livslang læring i Danmark}
\small
\renewcommand{\arraystretch}{1.3}
\begin{tabularx}{\linewidth}{@{} >{\leavevmode\bfseries\color{NCDNavy}\raggedright\arraybackslash}p{4.3cm} >{\raggedright\arraybackslash}X @{}}
\toprule
\multicolumn{2}{c}{\normalsize\bfseries\color{NCDRed} Emne 3B: Uddannelse med fokus på livslang læring i Danmark \quad\small\color{NCDGray}（仅含 stikord）} \\
\midrule
1. Dansk kontekst（终身学习文化） & livslang læring · efteruddannelse · AMU-kurser og VUC · aftenskoler · fagforening \\
\midrule
2. Fordele for individ \& job & følge med i AI og digitalisering · højere jobtryghed · personlig udvikling \\
\midrule
3. Ulemper / udfordringer & knap tid med fuldtidsarbejde og børn · træthed om aftenen · økonomi for selvstændige \\
\midrule
4. Egne erfaringer \& mål & danskuddannelse (DU3) · nye faglige kurser · sparring med kollegaer \\
\midrule
5. Fremtidsperspektiv（立场） & nødvendigt hele livet · støtte fra arbejdsgivere og fleksible onlinekurser \\
\bottomrule
\end{tabularx}

\vspace{0.45em}
\textbf{\color{NCDRed}Del 2 考官追问预测（Eksaminators opfølgende spørgsmål · 4 题）：}
\begin{enumerate}
  \item \textit{Vil du fortælle lidt mere om, hvorfor livslang læring er blevet så vigtig på det danske arbejdsmarked?}
  \item \textit{Kan du give nogle eksempler på, hvad du selv har lært som voksen i Danmark, og hvordan du har brugt det?}
  \item \textit{Hvad er fordelene og ulemperne ved at tage efteruddannelse online sammenlignet med fysisk undervisning?}
  \item \textit{Hvad synes du om mulighederne for voksen- og efteruddannelse i Danmark sammenlignet med i dit hjemland? Hvorfor?}
\end{enumerate}
\end{grammatikbox}

% ---------------------------------------------------------------------
\subsubsection*{Mundtlig Opgave 2 (Del 1 Individuel ca. 3 min. pr. kursist + Del 2 Paropgave ca. 4 min.)}

\begin{grammatikbox}{Opgaveark A \& B --- EMNE: Medier og digitalisering（个人图片情景提示卡）}
\small
\begin{tabularx}{\linewidth}{@{} X @{\hspace{0.4cm}} X @{}}
\begin{tcolorbox}[enhanced, colback=white, colframe=NCDTeal, boxrule=0.8pt, arc=3pt, left=7pt, right=7pt, top=6pt, bottom=6pt]
  \textbf{\color{NCDNavy}Opgaveark A：长者求助阅读 Digital Post}\\
  {\footnotesize\color{NCDGray}[图片 A：一位银发老人坐在明亮的餐桌旁，正用平板电脑打开 Digital Post 中的政府通知，身旁的一位年轻亲属或志愿者耐心地指着屏幕协助操作。]}
\end{tcolorbox}
&
\begin{tcolorbox}[enhanced, colback=white, colframe=NCDTeal, boxrule=0.8pt, arc=3pt, left=7pt, right=7pt, top=6pt, bottom=6pt]
  \textbf{\color{NCDNavy}Opgaveark B：咖啡馆里的低头族朋友}\\
  {\footnotesize\color{NCDGray}[图片 B：四位年轻朋友围坐在丹麦一家咖啡馆的木桌旁喝咖啡，但每个人都低着头专注地盯着各自的智能手机屏幕，彼此之间没有目光交流与交谈。]}
\end{tcolorbox}
\end{tabularx}
\end{grammatikbox}

\begin{emailbox}{TESTERARK --- 官方考官引导语、必问三题与双人辩论题卡（EMNE: Medier og digitalisering）}
\small\textcolor{NCDGray}{\textit{Instruktionen skal følges af tester. Hvis der vælges en anden sproglig formulering, skal instruktioner og spørgsmål indholdsmæssigt og sprogligt ligge så tæt på forlægget som muligt.}}

\vspace{0.35em}
\textbf{\color{NCDRed}Del 1 (individuel opgave -- ca. 3 min. pr. kursist)}\\
\textbf{\color{NCDNavy}KURSIST A}\\
Opgaveark A lægges foran kursist A, der præsenteres for emnet på følgende måde:
\begin{quote}
\textit{Emnet er 'Medier og digitalisering'. Her er et billede, som viser en ældre borger, der får hjælp til at læse sin Digital Post på en tablet.}
\end{quote}
Tester stiller derpå nedenstående obligatoriske spørgsmål:
\begin{enumerate}
  \item \textit{Der er nogle borgere, især ældre, der har brug for hjælp til at bruge offentlige digitale løsninger. Hvorfor tror du, det er sådan?}
  \item \textit{Hvilke fordele kan der være ved, at kontakten med det offentlige foregår digitalt?}
  \item \textit{Mener du, der kan være nogen ulemper ved, at næsten alt foregår digitalt? Hvorfor? / Hvorfor ikke?}
\end{enumerate}

\textbf{\color{NCDNavy}KURSIST B}\\
Opgaveark B lægges foran kursist B, der præsenteres for emnet på følgende måde:
\begin{quote}
\textit{Emnet er 'Medier og digitalisering'. Her er et billede, som viser en gruppe venner på en café, som alle kigger på deres smartphones.}
\end{quote}
Tester stiller derpå nedenstående obligatoriske spørgsmål:
\begin{enumerate}
  \item \textit{Der er mange mennesker, der bruger deres mobiltelefon hele tiden, også når de er sammen med venner. Hvorfor tror du, de gør det?}
  \item \textit{Hvilke fordele kan der være ved altid at være online på mobilen?}
  \item \textit{Mener du, der kan være nogen ulemper ved, at mobilen fylder så meget i sociale situationer? Hvorfor? / Hvorfor ikke?}
\end{enumerate}

\vspace{0.35em}
\hrule
\vspace{0.35em}
\textbf{\color{NCDRed}Del 2 (paropgave -- ca. 4 min.)}\\
Testen fortsætter som en samtale mellem de to kursister, hvor de udveksler synspunkter om et emne og begrunder disse. Tester præsenterer emnet med følgende input til samtalen:
\begin{quote}
\textit{I Danmark foregår næsten al kontakt med det offentlige i dag digitalt via MitID og Digital Post. Hvad synes I om det? I skal nu diskutere, hvilke fordele og ulemper, I mener, der kan være for den enkelte og for samfundet. I skal stille spørgsmål til hinanden og svare på hinandens spørgsmål.}
\end{quote}
\small\textcolor{NCDGray}{Tester fungerer som mediator og bryder ind, hvis samtalen ikke handler om emnet, eller hvis samtalen går i stå eller overtages helt af den ene kursist.}
\end{emailbox}

% =====================================================================
% DEL 4: RETTEARK, DETALJERET ANALYSE & HØJSCORE-MODELLER
% =====================================================================
\clearpage
\ncdsubsec{Del 4: Retteark, Detaljeret Analyse \& Højscore-Modeller}

\subsubsection*{4.1 官方 \textit{Retteark} 评分卡（29 point）与考前自查准则}

\small
\renewcommand{\arraystretch}{1.26}
\begin{longtable}{@{\extracolsep{\fill}} >{\leavevmode\bfseries\color{NCDNavy}\raggedright\arraybackslash}p{3.1cm} >{\raggedright\arraybackslash}p{8.6cm} >{\centering\arraybackslash}p{1.6cm} >{\centering\arraybackslash}p{2.2cm} @{}}
\toprule
\textbf{考核题型} & \textbf{Rettenøgle（标准答案与官方可接受变体）} & \textbf{分值} & \textbf{自测 / 评估} \\
\midrule
\endfirsthead
\toprule
\textbf{考核题型（续表）} & \textbf{Rettenøgle（标准答案与官方可接受变体）} & \textbf{分值} & \textbf{自测 / 评估} \\
\midrule
\endhead
\bottomrule
\endfoot
\textbf{Læsning Opgave 1}\newline（8 题） &
1. \textbf{B} (\textit{følelser}) \quad 2. \textbf{A} (\textit{Derfor}) \quad 3. \textbf{C} (\textit{bag}) \quad 4. \textbf{D} (\textit{dokumenterede})\newline
5. \textbf{A} (\textit{Alligevel}) \quad 6. \textbf{B} (\textit{skærpe}) \quad 7. \textbf{C} (\textit{skeptisk over for}) \quad 8. \textbf{D} (\textit{I sidste ende})\newline
\textcolor{NCDGray}{\small （例题 0 为 C \textit{fængende}）} & 8 point & \rule{0.65cm}{0.4pt} / 8 \\
\midrule
\textbf{Læsning Opgave 2}\newline（5 题） &
1. \textbf{F} \quad 2. \textbf{C} \quad 3. \textbf{H} \quad 4. \textbf{E} \quad 5. \textbf{D}\newline
\textcolor{NCDGray}{\small （例句 0 为 A｜多余干扰项：B, G）} & 5 point & \rule{0.65cm}{0.4pt} / 5 \\
\midrule
\textbf{Læsning Opgave 3}\newline（10 题 · 无词库） &
1. \textbf{hvor} \quad 2. \textbf{ikke} \quad 3. \textbf{med} \quad 4. \textbf{som} eller \textbf{der} \quad 5. \textbf{ved}\newline
6. \textbf{der} eller \textbf{som} \quad 7. \textbf{Desuden} (eller \textbf{Derudover} / \textbf{Samtidig}) \quad 8. \textbf{til}\newline
9. \textbf{så} \quad 10. \textbf{om} & 10 point & \rule{0.65cm}{0.4pt} / 10 \\
\midrule
\textbf{Læsning Opgave 4}\newline（6 题） &
1. \textbf{B} \quad 2. \textbf{C} \quad 3. \textbf{A} \quad 4. \textbf{B} \quad 5. \textbf{C} \quad 6. \textbf{A} & 6 point & \rule{0.65cm}{0.4pt} / 6 \\
\midrule
\textbf{\color{NCDRed}Læsning 总计} & \textbf{Opgave 1--4 合订总评（50 分钟 · 禁用词典）} & \textbf{29 point} & \rule{0.65cm}{0.4pt} / 29 \\
\midrule
\textbf{Skrivning Opgave 1}\newline（私人邮件回信） &
词数严格 \textbf{$\ge$ 125 ord}；全面回应 Jonas 全部 3 个维度（关闭晚间新闻推送的睡眠/信息焦虑体验、周末纸媒订阅 vs 免费新闻播客选择与理由、父母社媒晒娃隐私立场）；自然融入情态小品词。 & 整体评估 & Bestået / Ikke bestået \\
\midrule
\textbf{Skrivning Opgave 2}\newline（观点辩论帖） &
词数严格 \textbf{$\ge$ 100 ord}；清晰表明是否同意 Søren Holm“社交媒体设 16 岁法定年龄门槛”观点，运用让步反驳并兼顾青少年心理健康与监管可行性提出方案。 & 整体评估 & Bestået / Ikke bestået \\
\midrule
\textbf{Mundtlig Opgave 1}\newline（Emneopgave） &
仅带纯 \textit{stikord} 提示卡；Del 1 陈述 1--2 分钟兼顾丹麦媒体/终身学习体制与个人经验；Del 2 以较长话轮回答追问并明确论证利弊与原因。 & 整体评估 & Bestået / Ikke bestået \\
\midrule
\textbf{Mundtlig Opgave 2}\newline（Ind. + Paropgave） &
Del 1（约 3 分钟/人）：回答图片三问（原因 / 优点 / 缺点与理由）；Del 2（约 4 分钟）：双人互动互问互答，必须兼顾\textbf{对个人与对社会}两方面利弊。 & 整体评估 & Bestået / Ikke bestået \\
\end{longtable}
\normalsize

\begin{tipbox}{考前自查准则与失分诊断（Self-Check Checklist）}
\small
\begin{itemize}
  \item \textbf{总分达标线参考}：官方 \textit{Retteark} 仅记 29 分原始总分，各语言中心依 SIRI 指引划定的及格线通常约为 60--65\%（即 \textbf{18--19 分以上}）。建议冲刺自测稳定拿到 \textbf{20+/29}。
  \item \textbf{Opgave 3 无词库与介词易错点排查}：若在 Opgave 3 低于 7/10，请重点检查三类失分：\ding{172} 极易混淆的固定介词（\textit{overbevist \textbf{om}} 切勿受 \textit{sikker på} 干扰误填 \textit{på}！另有 \textit{svært \textbf{ved}}, \textit{penge \textbf{til}}）；\ding{173} 关系代词与副词（\textit{i en tid, \textbf{hvor}} 表时间背景；\textit{emner, \textbf{der/som}} 作主语）；\ding{174} 比较级对称结构（\textit{lige så meget ... \textbf{som} ...}）与结果连词 \textit{\textbf{så}}。
  \item \textbf{本卷语言靶点与考场红线}：\ding{172} 写作中主动交替使用因果连接（\textit{som følge heraf, følgelig, hvorved}）与传闻推测情态（\textit{vist, nok, angiveligt}），避免全篇只用 \textit{fordi / så}；\ding{173} 写完务必逐词点数（Opgave 1 $\ge$ 125 词、Opgave 2 $\ge$ 100 词）；\ding{174} 口语 \textit{Emneopgave} 绝不携带完整句子，Opgave 2 Del 2 务必覆盖 \textit{for den enkelte} 与 \textit{for samfundet} 双维度。
\end{itemize}
\end{tipbox}

% ---------------------------------------------------------------------
\subsubsection*{4.2 阅读 Opgave 1--4 逐题解析、全文中文精译与高频词表}

\begin{answerbox}
\textbf{\color{NCDNavy}【一、Opgave 1（Nyheder på sociale medier og kildekritik）逐题语法与语义解析】}
\begin{itemize}
  \item \textbf{(0) C (\textit{fængende}) [Eksempel]}：与 \textit{korte} 并列修饰 \textit{videoer}，说明社媒新闻常以简短且“抓人眼球/有吸引力的（\textit{fængende}）”短视频形式呈现。
  \item \textbf{(1) B (\textit{følelser})}：后文直接以 \textit{som vrede, frygt eller begejstring}（如愤怒、恐惧或激动）举例，这些均属于强烈“情绪/情感（\textit{følelser}）”。A（规则）、C（习惯）、D（支出）无法搭配该举例。
  \item \textbf{(2) A (\textit{Derfor})}：前句阐明“互动越多，传播越快”，本句顺承结果：“\textbf{因此}，戏剧化且片面的论断往往比冷静平衡的分析更常被推送”，位于句首触发 V2 倒装（\textit{Derfor bliver...}）。\textbf{排雷}：B（相反）、C（相对而言）、D（尽管如此）均为转折方向，与因果逻辑背道而驰。
  \item \textbf{(3) C (\textit{bag})}：固定动词短语 \textit{stå bag noget}（在某事背后操盘/负责某条信息），追问网上帖子究竟由谁发布。
  \item \textbf{(4) D (\textit{dokumenterede})}：编辑部与专业记者过去担保信息在见报前均经过“查证落实/有据可依（\textit{dokumenterede}）”。\textbf{排雷}：A（编造的）、B（片面的）、C（随机的）均属贬义，与 \textit{garant for}（担保）矛盾。
  \item \textbf{(5) A (\textit{Alligevel})}：上一段揭示匿名账号与假新闻乱象，本段话锋一转：“\textbf{尽管如此}，解决办法未必是完全避开社交媒体”。B（因此）、C（遗憾的是）、D（以前）均无法承接“先抑后扬”的转折。
  \item \textbf{(6) B (\textit{skærpe})}：标准媒体评论搭配 \textit{skærpe sit kritiske blik}（擦亮/敏锐自己的批判视角）。A（忘记）、C（隐藏）、D（放弃）与点击转发前核查完全相反。
  \item \textbf{(7) C (\textit{skeptisk over for})}：当夸张标题缺乏日期或具体证据时，应当对内容保持怀疑态度（\textit{være skeptisk over for indholdet}）。A（满意）、B（依赖）、D（喜欢）均不合风控语境。
  \item \textbf{(8) D (\textit{I sidste ende})}：全篇收尾金句，“归根结底/说到底（\textit{I sidste ende}），未来的新闻习惯不是抛弃技术，而是向屏幕所见内容提出正确的问题”。A（相反）、B（例如）、C（遗憾的是）均不能总括核心主旨。
\end{itemize}
\end{answerbox}

\begin{transbox}
\small
\textbf{\color{NCDTeal}[Opgave 1 全文中文精译]}\\
如今越来越多丹麦人在 Instagram 和 TikTok 等社交媒体上阅读每日新闻，而不是购买报纸。这样做既方便又免费，因为手机总在口袋里，而且新闻以简短抓人的视频形式呈现。然而，快速刷新的信息流也带来了强烈的信源批判需求。\\[3pt]
社交媒体的一大挑战在于，算法首先会推送能激发愤怒、恐惧或激动等强烈情绪的内容。一条帖子的评论与转发越多，就会越快地扩散给成千上万的用户。因此，戏剧化且片面的说法往往比传统主流媒体冷静、多维度的分析更容易频繁出现。\\[3pt]
此外，网上某条帖子的背后究竟是谁，也并不总是清晰明确。过去，由受过专业训练的记者组成的编辑部会担保信息在见报前经过充分核实与佐证；而如今，任何人都能创建一个看起来像正规新闻媒体的精美账号，传播据称来自可信来源的谣言或篡改图像。\\[3pt]
尽管如此，解决办法未必是彻底远离社交媒体。研究人员指出，人们只需养成习惯，在点击“分享”前花几秒钟擦亮批判的眼光。例如，可以查证是否有多家独立新闻媒体报道同一事件，并留意究竟是谁以专家身份发声。如果一条耸人听闻的标题缺少日期或具体证据，就应当对内容保持警惕和怀疑。\\[3pt]
目前，多所丹麦公立小学和图书馆正面向儿童和青少年开展数字素养与事实核查教育，以帮助下一代具备在信息洪流中辨别方向的能力。归根结底，未来的新闻阅读习惯并不在于排斥科技，而在于学会对屏幕上呈现的内容提出恰当的问题。
\end{transbox}

\begin{answerbox}
\textbf{\color{NCDNavy}【二、Opgave 2（Avisen der gik tilbage til papir）逐题衔接解析】}
\begin{itemize}
  \item \textbf{(0) A [Eksempel]}：与前两句共同交代五年前放弃纸媒的背景：印厂与派送太贵，“与此同时越来越多年轻市民在公交上或午休时用手机刷新闻”，自然衬托下文 \textit{Men allerede efter to år...} 的反转。
  \item \textbf{(1) F}：前句写网站虽有点击但阅读时长暴跌（读者点开 Facebook 看十秒就划走），\textbf{F} 用因果连词 \textit{Derfor} 顺承：“因此付费订阅人数下降，因为人们不愿为转眼就忘的短文付费”，随即接 \textit{Samtidig}（老订户写信说怀念周六清晨看实体报）。
  \item \textbf{(2) C}：前句宣布重推双周厚实纸质版，主打报道与本地人物特写，\textbf{C} 阐明编辑部定位：“计划从来不是在速度上与互联网竞争，而是创造一片供深度沉浸阅读的喘息空间”，紧接着下一句以 \textit{I stedet for at jage hurtige nyheder...}（不再追逐车祸快讯）精准承接。
  \item \textbf{(3) H}：前句指印刷纸价暴涨、预算脆弱，需要另辟蹊径，\textbf{H} 给出行动：“因此他们与本地商铺达成紧密合作，涵盖广告投放与收银台售报”，而下句代词 \textit{Det}（“这”）直接回指该项合作大获成功，因为面包店、书店与文化馆都乐意支持广告能被从容细读的本地媒体。
  \item \textbf{(4) E}：前句描述纸质报整个周末放在餐桌上，大人和孩子都会翻阅，\textbf{E} 提炼纸媒物理优势：“与锁屏手机不同，一张报纸可以在几天内由家庭成员实体传阅共享”，从而直接导出下句的 \textit{På den måde}（正是通过这种方式，引发餐桌对话）。
  \item \textbf{(5) D}：前句提出“纸媒与网络并不排斥，反而能共同强化本地新闻”，\textbf{D} 精准对比两者分工：“当网络在平日提供快捷更新时，纸质版则为本地社区留出了思考与信任的时间”，自然引出主编最后“好新闻需要时间”的总结引语。
  \item \textbf{多余干扰项排雷}：\textbf{B} 称“绝大多数丹麦人现在只读英文国际报纸”，与本文本地社区报主题完全冲突；\textbf{G} 称“市政府为省电费关闭公共图书馆”，纯属借“省钱/本地”词汇拼凑的离题干扰句。
\end{itemize}
\end{answerbox}

\begin{transbox}
\small
\textbf{\color{NCDTeal}[Opgave 2 全文中文精译]}\\
五年前，当日德兰一座沿海小城的本地报纸《峡湾邮报》决定停办纸质版、转为纯数字化出版时，这听起来像是唯一明智的出路。印刷厂和投递成本相对于广告收入已变得过于高昂。\textbf{（0-A）与此同时，越来越多的年轻市民习惯在公交车上或午休时用智能手机看新闻。}但仅仅过了两年，主编梅特·克罗格和她的记者们便开始察觉到一种意想不到的趋势，让他们对当初的决定产生了怀疑。\\[3pt]
尽管网站依然有不少点击，阅读时长却断崖式下跌。许多读者只是通过脸书点开文章看十秒钟，便继续向下滑动。\textbf{（1-F）因此，付费订阅者的数量不断下滑，因为人们不愿为看过转瞬即忘的短文掏钱。}与此同时，老订户们写信表示，他们非常怀念周六上午在早餐桌旁翻看实体报纸的时光——没有电子屏幕，也没有闪烁的广告。\\[3pt]
经过长时间思考，编辑部迈出了被不少人称为“老派”的一步。他们决定将《峡湾邮报》重新改版为双周发行的厚实纸质报，选用质感厚重的纸张，刊载更多本地深度报道与人物专访。\textbf{（2-C）他们的计划从来不是在速度上与互联网竞争，而是为深度阅读营造一个喘息空间。}该报不再追逐交通事故之类的快讯，而是把重心放在背景深挖文章、观点评论以及本地热心公益人物的故事上。\\[3pt]
尽管如此，重启初期在财务上依然十分脆弱，因为印刷费和纸张成本大幅上涨。编辑部必须发挥创意，在不牺牲质量的前提下实现收支平衡。\textbf{（3-H）于是，他们与本地商铺在广告投放和收银台代售报纸两方面建立了紧密合作。}幸运的是，这一举措被证明非常成功，因为镇上的面包房、书店和文化中心都乐于支持一家能让自家广告被市民从容细读的媒体。\\[3pt]
如今，《峡湾邮报》收获了数百名新订户，其中出人意料地包括许多三十多岁的有孩家庭。不少读者表示，报纸整个周末都会摊放在餐桌上，大人和孩子都会随手翻看。\textbf{（4-E）与锁屏的手机不同，一份报纸可以在几天时间里在家庭成员之间实体传阅。}正是以这种方式，纸质报在晚餐桌上激发了关于城里学校和本地社团的话题讨论——而当每个人都各自低头看屏幕时，这几乎很少发生。\\[3pt]
不过，主编梅特·克罗格并不认为所有丹麦报纸都应该关停数字平台。因为纸媒与互联网并不相互排斥，反而能携手增强本地新闻的力量。\textbf{（5-D）当网络在工作日提供快速更新时，纸质版则为本地社区保留了沉淀思考与建立信任的时间。}“对我们来说，回归纸媒并不是倒退一步，而是一次提醒：好新闻需要时间——无论是对撰写它的人，还是对阅读它的人而言，”她最后总结道。
\end{transbox}

\begin{answerbox}
\textbf{\color{NCDNavy}【三、Opgave 3（Nadia arbejder som børnebibliotekar）无词库填词解析】}
\begin{itemize}
  \item \textbf{(1) hvor}：时间关系副词，承接 \textit{i en tid, hvor...}（在一个……的时代），相当于“在这个时代里”，为丹麦语标准书面结构。
  \item \textbf{(2) ikke}：固定加强否定副词搭配 \textit{slet ikke}（根本不，一点也不），置于从句主语 \textit{børn} 之后、情态动词 \textit{gider} 之前。
  \item \textbf{(3) med}：固定介词短语 \textit{sammen med nogen}（和某人一起）。
  \item \textbf{(4) som}：比较对称结构 \textit{lige så meget om A som om B}（既关乎 A 同样也关乎 B），必须使用 \textit{som}。
  \item \textbf{(5) ved}：固定搭配 \textit{have svært ved at + 不定式}（做某事有困难）。
  \item \textbf{(6) der / som}：定语从句关系代词修饰先行词 \textit{emner} 并在从句中作主语（\textit{...interesserer barnet}），故 \textit{der} 与 \textit{som} 均完全正确。
  \item \textbf{(7) Desuden / Derudover / Samtidig}：句首递进/并列副词，衔接“此外，她还在周末组织小型读书会和创意工坊”，注意首字母须大写且紧随其后发生 V2 倒装（\textit{arrangerer hun...}）。
  \item \textbf{(8) til}：固定搭配 \textit{penge til at købe...}（用于购买……的经费/钱），表用途不可误填 \textit{for}。
  \item \textbf{(9) så}：结果从句连词（“从而使暑假前繁忙的时期变得更容易应付”）；后面接完整主谓结构且限于一词，不可用 \textit{for}。
  \item \textbf{(10) om}：核心高频陷阱！形容词 \textit{overbevist} 固定搭配介词 \textit{om}（\textit{være overbevist om, at...} 确信……），切忌受 \textit{sikker på} 干扰误填 \textit{på}！
\end{itemize}
\end{answerbox}

\begin{transbox}
\small
\textbf{\color{NCDTeal}[Opgave 3 全文中文精译]}\\
纳迪亚在奥胡斯担任儿童图书管理员已有七年。她最重要的任务，是在电子屏幕比以往任何时候都更加占据生活的时代，唤醒孩子们对故事的热爱。尽管许多成年人以为孩子如今根本不再愿意翻开实体书，但纳迪亚每周都在借阅台前看到长长的队伍。\\[3pt]
上午时分，图书馆经常迎来幼儿园和小学班级的来访。纳迪亚和孩子们一起坐在地板上的软垫上，大声朗读童话故事和新绘本。之后大家一起讨论图画并互相提问。对纳迪亚来说，这份工作关乎陪伴的程度，丝毫不亚于书本本身。\\[3pt]
除了朗读活动，纳迪亚还帮助家长挑选适合孩子年龄的图书。有些孩子很难开始阅读，于是她会找来有趣的漫画或关于孩子感兴趣主题的书。此外，她还在周末组织小型读书俱乐部和工坊，让家庭可以免费相聚。\\[3pt]
这份工作最困难的地方在于，图书馆的预算已多次遭到削减。例如，常常缺乏资金购买足够的新儿童读物。尽管如此，当策划各类活动时，同事们依然很善于彼此扶持，从而让暑假前繁忙的时期变得更容易应对。\\[3pt]
每当纳迪亚看到一个孩子怀里抱着一摞书开开心心地回家时，她就非常清楚自己为何热爱这份工作。她坚信，强大的儿童图书馆能给予所有孩子平等的机会——完全不受父母钱包厚薄的限制。
\end{transbox}

\begin{answerbox}
\textbf{\color{NCDNavy}【四、Opgave 4（6 则实用文本）定位与排雷解析】}
\begin{itemize}
  \item \textbf{1. B}：原文末句明确写 \textit{Bliver du fritaget, får du fremover offentlige breve sendt med fysisk post}，与 B 完全同义。\textbf{排雷}：A 虚构“65 岁自动豁免”（原文须具合理事由并亲自或委托前往 Borgerservice 申请）；C 与“亲自到场或委托他人持授权书前往 Borgerservice”相反。
  \item \textbf{2. C}：原文末句 \textit{Derfor bør du altid vente med at lukke din gamle bankkonto, indtil din nye NemKonto er aktiv} 与 C 精准对应。\textbf{排雷}：A 缩小范围，原文举例包括退税、SU、工资与住房补贴；B 否认可在 nemkonto.dk 自行用 MitID 修改或找新银行代办。
  \item \textbf{3. A}：短信写明讲座改至下周二 19 点，且在周一 12 点前回复即可全额退票，对应 A。\textbf{排雷}：B 违背 \textit{Din billet gælder automatisk}；C 颠倒退款截止时间（须在周一 12 点前而非周二讲座后）。
  \item \textbf{4. B}：原文指出拒绝统计与营销类 Cookie 后，用户仍有权浏览网站常规内容（\textit{har du stadig ret til at læse hjemmesidens almindelige indhold}），对应 B。\textbf{排雷}：A 忽略“技术必要型 Cookie 无需同意”；C 违背“可随时更改选择”。
  \item \textbf{5. C}：Emilie 自己用完了教材，想从明天下午起借给 Tobias，对应 C。\textbf{排雷}：A 颠倒借出与借入双方；B 错把“若你还没买书”曲解为“让 Tobias 给她买新书”。
  \item \textbf{6. A}：原文建议外出办理网银或 MitID 时不要用未知公共 Wi-Fi，改用手机自带移动数据或安全 VPN，对应 A。\textbf{排雷}：B 与“连接往往未加密（\textit{ofte ikke krypteret}）”完全相反；C 把“审批支付时用 MitID”张冠李戴为“连任何 Wi-Fi 都要用 MitID”。
\end{itemize}
\end{answerbox}

\begin{transbox}
\small
\textbf{\color{NCDTeal}[Opgave 4 六则短文中文对照]}\\
\textbf{1. Digital Post 与豁免规定（borger.dk）：}丹麦所有年满 15 周岁的居民原则上通过 Digital Post 接收公共机构来信（可使用 MitID 在 borger.dk、e-Boks 或 mit.dk 阅读）。如因疾病、残障、语言困难或家中缺乏网络无法使用数字自助服务，可申请豁免。申请须本人前往市政市民服务中心（Borgerservice），或由亲属持书面授权书代办。获批豁免后，日后公共机构来信将改以实体纸质信件寄送。\\
\textbf{2. 在银行变更 NemKonto（nemkonto.dk）：}丹麦所有年满 18 周岁居民均须拥有一个 NemKonto，用于接收退税、SU、工资或住房补贴等公共款项。如更换银行或想改用其他账户，可使用 MitID 在 nemkonto.dk 自行变更，或请新银行代为登记。请注意变更生效可能需几个工作日，因此务必等新 NemKonto 激活后再关闭旧银行账户，以免款项延误。\\
\textbf{3. 文化中心晚间讲座改期短信：}嗨阿斯特丽德，周四关于假新闻的晚间讲座主讲人遗憾生病了。因此讲座改期至下周二 19:00 在同一礼堂举行。你的门票对新日期自动有效；但若当天无法前来，请于周一中午 12 点前回复本短信，即可获得全额退款。祝好，韦斯特布罗文化中心。\\
\textbf{4. Cookie 与网络数据保护（datatilsynet.dk）：}访问丹麦网站时，未经你明确同意，网站不得收集你的行为数据用于统计或定向营销。网站必须提供真实选择，使“拒绝”与“接受”同样容易。仅维持购物车或登录运行的技术必要型 Cookie 可免同意使用。若拒绝统计与营销 Cookie，你仍有权阅读网站常规内容，并可随时撤回或更改决定。\\
\textbf{5. 同学借教材短信：}嗨托比亚斯，你下周写作业还需要咱们那本社会学教材吗？我自己的项目已经做完了，如果你还没买书的话，非常欢迎从明天下午起借我的去用。只要你今晚回复我一下，我明天下午 2 点可以带到食堂给你。祝好，埃米莉。\\
\textbf{6. 公共 Wi-Fi 安全与 MitID（sikkerdigital.dk）：}在咖啡馆、火车或机场使用开放式公共 Wi-Fi 虽然方便，但连接往往并未加密。因此不建议在未知开放网络下登录网银、查看敏感医疗数据或使用 MitID 确认付款。在外处理重要数字事务时，请改用手机移动数据或开启安全 VPN 连接。同时，滑动 MitID 确认前务必仔细核对屏幕文字，遇到非本人发起的请求请立即拒绝。
\end{transbox}

\paragraph*{核心高频词表（Gloser｜15 个 B1+/B2- 完整屈折词条）}

\begin{glosertable}
名词 & en kildekritik & \textit{(kildekritikken, \textnormal{通常无复数})} & 信源批判，信息甄别能力 \\
名词 & en algoritme & \textit{(algoritmen, algoritmer, algoritmerne)} & 算法 \\
名词 & en påstand & \textit{(påstanden, påstande, påstandene)} & 论断，说法，主张 \\
形容词 & troværdig & \textit{(troværdig, troværdigt, troværdige)} & 可信的，可靠的 \\
动词 [har] & at efterprøve & \textit{(efterprøver, efterprøvede, har efterprøvet)} & 核实，验证，检验 \\
动词 [har] & at dokumentere & \textit{(dokumenterer, dokumenterede, har dokumenteret)} & 证实，以证据说明 \\
名词 & en fordybelse & \textit{(fordybelsen, \textnormal{通常无复数})} & 沉浸，深入钻研，专注 \\
名词 & en abonnent & \textit{(abonnenten, abonnenter, abonnenterne)} & 订阅者，订户 \\
名词 & et pusterum & \textit{(pusterummet, pusterum, pusterummene)} & 喘息空间，静心时刻（单复同形） \\
动词 [har] & at fritage & \textit{(fritager, fritog, har fritaget)} & 免除，豁免（不规则！） \\
名词 & et samtykke & \textit{(samtykket, samtykker, samtykkerne)} & 同意，授权，许可 \\
动词 [har] & at omgå & \textit{(omgår, omgik, har omgået)} & 绕过，规避（不规则！） \\
名词 & en aldersgrænse & \textit{(aldersgrænsen, aldersgrænser, -erne)} & 年龄限制，年龄门槛 \\
名词 & en efteruddannelse & \textit{(efteruddannelsen, efteruddannelser, -erne)} & 继续教育，在职进修 \\
名词 & et privatliv & \textit{(privatlivet, \textnormal{通常无复数})} & 隐私，私人生活 \\
\end{glosertable}

% ---------------------------------------------------------------------
\subsubsection*{4.3 写作双档高分示范（Opgave 1 $\ge$ 125 ord \& Opgave 2 $\ge$ 100 ord）}

\begin{emailbox}{Skrivning Opgave 1 范文一：B1+ 稳过版（154 ord｜三问全覆 · 情态小品词自然落地）}
Kære Jonas

\vspace{0.25em}
Tusind tak for din mail! Det er \da{jo} altid så hyggeligt at høre fra dig.

\vspace{0.25em}
Det går faktisk rigtig godt uden nyhedsnotifikationer om aftenen. I starten var jeg lidt bange for at overse noget vigtigt, men efter få dage mærkede jeg, at jeg sover langt roligere. Man behøver \da{jo} ikke følge med i alt hele tiden.

\vspace{0.25em}
Jeg forstår godt, at du vil ændre dine egne nyhedsvaner. \dabold{Hvis jeg var dig, ville jeg vælge} at abonnere på en trykt weekendavis. En podcast er fin i bussen, men en fysisk avis giver en helt særlig ro lørdag morgen uden forstyrrelser fra telefonen.

\vspace{0.25em}
Hvad angår din søster, synes jeg ikke, at forældre bør dele billeder af små børn på sociale medier. Børn kan \da{jo} ikke selv sige ja eller nej, og billederne kan blive liggende på nettet i mange år. Derfor bør man passe på deres privatliv.

\vspace{0.25em}
Skriv snart, hvad du vælger!

\vspace{0.25em}
Kærlig hilsen\\
Kasper

\vspace{0.3em}
\rule{\linewidth}{0.4pt}\\[0.2em]
\small\textbf{\color{NCDTeal}[考官点评 · 154 ord]}：稳居 145--165 词安全区。逐一回答 Jonas 全部三个问题（晚间关推送的睡眠改善、选周末纸媒的建议与理由、反对父母晒娃的隐私顾虑），恰当嵌入 \textit{jo} 拉近朋友口吻，并规范运用第二条件句（\textit{Hvis jeg var dig, ville jeg...}）作建议。
\end{emailbox}

\begin{emailbox}{Skrivning Opgave 1 范文二：B2- 高分冲刺版（182 ord｜因果连接 + 传闻情态 + 让步 + 被动语态）}
Kære Jonas

\vspace{0.25em}
Hvor er det dejligt at høre fra dig igen! Det er \da{vist} allerede et par måneder siden, vi sidst sås.

\vspace{0.25em}
At slå notifikationerne fra har \da{egentlig} gjort en overraskende stor forskel for mig. Den konstante strøm af dramatiske overskrifter stresser kroppen, \dabold{hvorved} man mister dyb nattesøvn. I dag tjekker jeg kun udvalgte medier to gange dagligt, og \dabold{som følge heraf} sover jeg langt bedre om natten uden at savne noget væsentligt.

\vspace{0.25em}
Til dit eget nyhedsforbrug vil jeg klart anbefale en trykt weekendavis frem for daglige podcasts. \dabold{Ganske vist} er podcasts gratis og praktiske under transport, \dabold{men} avisen lægger i højere grad op til fordybelse og grundig kildekritik. Desuden har lokale aviser \dabold{nok} ekstra meget brug for abonnenter netop nu.

\vspace{0.25em}
Angående din søster mener jeg, at børns digitale privatliv \dabold{bør respekteres}. Billeder, der \dabold{deles} offentligt, kan \dabold{angiveligt} let kopieres og gemmes af fremmede uden forældrenes vidende. \dabold{Følgelig} er det langt klogere at sende familiebilleder i lukkede beskeder til bedsteforældrene i stedet for at lægge dem åbent ud på nettet.

\vspace{0.25em}
Lad os snart mødes til en god kop kaffe!

\vspace{0.25em}
Bedste hilsner\\
Kasper

\vspace{0.3em}
\rule{\linewidth}{0.4pt}\\[0.2em]
\small\textbf{\color{NCDNavy}[高分亮点 · 182 ord]}：\ding{172} 精准嵌入本卷因果三剑客 \textit{hvorved / som følge heraf / Følgelig} 与传闻推测情态 \textit{vist / nok / angiveligt}；\ding{173} 运用 \textit{Ganske vist..., men...} 对比播客与周末纸媒深度；\ding{174} 规范使用 \textit{-s} 被动（\textit{bør respekteres, deles}），语域地道老练。
\end{emailbox}

\begin{emailbox}{Skrivning Opgave 2 范文一：B1+ 稳过版（正文 133 ord / 含标题 140 ord｜正反双面权衡）}
\textbf{\color{NCDNavy}Beskyt børn bedre uden et totalt forbud}

\vspace{0.25em}
Søren Holm tager et vigtigt problem op i sit indlæg, men jeg er kun delvist enig i hans forslag om en streng aldersgrænse på 16 år for sociale medier.

\vspace{0.25em}
På den ene side har han ret i, at sociale medier kan skade børns søvn, selvtillid og koncentration. Mange unge bruger alt for meget tid på korte videoer og oplever desværre også mobning på nettet.

\vspace{0.25em}
På den anden side er et totalt forbud indtil 16-årsalderen svært at kontrollere i praksis. Unge finder ofte måder at snyde på, og samtidig bruger mange teenagere også sociale medier til at snakke med venner og lære om spændende fritidsinteresser.

\vspace{0.25em}
Derfor bør vi hellere stille strengere krav til tech-giganterne og undervise i digital dannelse i skolen. Så lærer de unge at bruge nettet fornuftigt, før de bliver 16 år.

\vspace{0.3em}
\rule{\linewidth}{0.4pt}\\[0.2em]
\small\textbf{\color{NCDTeal}[考官点评 · 133/140 ord]}：紧扣 B1+ 稳过区间（125--145 ord）。开篇表明“部分同意但质疑一刀切年龄禁令”，中段采用 \textit{På den ene side / På den anden side} 客观权衡睡眠压力与社交益处及监管漏洞，末段提出“规制平台+学校数字素养教育”的务实方案。
\end{emailbox}

\begin{emailbox}{Skrivning Opgave 2 范文二：B2- 高分冲刺版（正文 159 ord / 含标题 166 ord｜让步反驳 + 因果 + 推测情态 + 被动）}
\textbf{\color{NCDNavy}Regulering af tech-giganter slår en rigid aldersgrænse}

\vspace{0.25em}
Søren Holm rejser en yderst relevant bekymring for unges trivsel, men jeg er grundlæggende skeptisk over for, at en lovfæstet aldersgrænse på 16 år alene løser problemet.

\vspace{0.25em}
\dabold{Ganske vist} skaber afhængighedsskabende algoritmer på TikTok og Instagram mistrivsel blandt børn, og et øget skærmforbrug har \dabold{angiveligt} svækket både unges nattesøvn og koncentrationsevne. \dabold{Ikke desto mindre} kan et ensidigt forbud let omgås med falske fødselsdatoer, \dabold{hvorved} mange unge blot flytter over på langt mere skjulte og uregulerede platforme.

\vspace{0.25em}
\dabold{Hvis alle under 16 år blev udelukket} fra sociale medier, \dabold{ville} mange unge \dabold{nok} også miste vigtige fællesskaber om fritidsinteresser og politisk debat. \dabold{Som følge heraf} risikerer de at stå helt uforberedte, den dag forbuddet pludselig ophæves.

\vspace{0.25em}
Efter min mening bør løsningen rettes mod platformenes design frem for kun mod brugerne. Tech-virksomheder \dabold{bør forpligtes} til at fjerne skadelige algoritmer for mindreårige og indføre klare tidsgrænser. \dabold{Følgelig} \dabold{beskyttes} børn og unge effektivt, mens forældre og skoler samtidig støtter deres sunde digitale dannelse i hverdagen.

\vspace{0.3em}
\rule{\linewidth}{0.4pt}\\[0.2em]
\small\textbf{\color{NCDNavy}[高分亮点 · 159/166 ord]}：\ding{172} 严格落于 155--180 词高分带，开篇使用 \textit{skeptisk over for} 与 \textit{Ganske vist..., Ikke desto mindre...}；\ding{173} 密嵌因果连接 \textit{hvorved / Som følge heraf / Følgelig} 与传闻推测 \textit{angiveligt / nok}；\ding{174} 假拟条件嵌套 \textit{blive} 被动（\textit{Hvis alle under 16 år blev udelukket..., ville...}）与规范 \textit{-s} 被动（\textit{bør forpligtes, beskyttes}）精准收尾。
\end{emailbox}

% ---------------------------------------------------------------------
\subsubsection*{4.4 口语示范逐字稿、核心 IPA 与双人社会议题辩论}

\begin{transbox}
\small
\textbf{\color{NCDNavy}[Emneopgave 3A 口语陈述示范：Moderne livsstil: nyhedsvaner og kildekritik｜162 ord · 约 1.5 分钟]}\\
\textit{Jeg vil gerne fortælle om moderne livsstil i Danmark med fokus på nyhedsvaner og kildekritik \ipa{ˌkʰiləkʰʁiˈtˢiɡ}. I Danmark følger langt de fleste borgere med i samfundet gennem etablerede public service-medier som DR og TV 2, men stadig flere unge henter primært deres nyheder på sociale medier. Fordelen ved den digitale nyhedsstrøm er selvfølgelig, at man får et hurtigt og gratis overblik \ipa{ˈʌu̯ɐˌbleɡ} på mobilen døgnet rundt. Ulempen er imidlertid, at algoritmer ofte favoriserer voldsomme overskrifter frem for grundig journalistik, hvorved falske nyheder og ekkokamre \ipa{ˈɛɡoˌkʰɑmʁɐ} let spredes. I min egen hverdag prøver jeg derfor at være meget bevidst om mine daglige nyhedsvaner: når jeg ser en dramatisk historie på nettet, sammenligner jeg altid mindst to troværdige \ipa{tˢʁoˈvɛːˀɐðiə} og uafhængige kilder, før jeg tror på indholdet. Desuden har jeg slået push-beskeder fra om aftenen og læser i stedet længere baggrundsartikler i en trykt avis i weekenden. Samlet set mener jeg, at stærk digital dannelse allerede i folkeskolen er helt afgørende, så alle lærer at gennemskue manipuleret information.}\\[4pt]
\textbf{\color{NCDTeal}[参考译文]}：我想聊聊丹麦的现代生活方式，重点是新闻习惯与信源批判。在丹麦，绝大多数市民通过 DR 和 TV 2 等主流公共服务媒体了解社会动态，但越来越多年轻人主要从社交媒体获取新闻。数字新闻流的优点当然是能全天候在手机上获得快捷且免费的概览。然而缺点在于，算法往往偏爱耸动标题胜过扎实新闻报道，从而导致假新闻与信息茧房轻易扩散。因此在我自己的日常生活中，我尽量对每日新闻习惯保持高度自觉：每当在网上看到戏剧性的消息，我在相信其内容前总会对比至少两个可信且独立的信源；此外我在晚上关闭了推送通知，改为周末在纸质报上阅读长篇深度报道。总体而言，我认为从小学起开展扎实的数字素养教育至关重要，能让每个人都学会识破被操纵的信息。
\end{transbox}

\paragraph*{Emneopgave 3A Del 2 考官追问与高分应答示范（4 题）：}
\dialogturn{1}{Eksaminator}{Vil du fortælle lidt mere om, hvordan du selv vurderer, om en nyhed på internettet er troværdig?}{troværdig \ipa{tˢʁoˈvɛːˀɐði} · vurderer \ipa{vuˈdeːˀɐ}}{你能再多讲讲自己是如何判断网上一条新闻是否可信的吗？}
\dialogturn{2}{Kursist}{Ja, først kigger jeg på, hvem der står bag artiklen eller opslaget, og om der nævnes konkrete eksperter eller myndigheder med navn. Hvis kilden er anonym, eller hvis tonen virker meget vred og ensidig, tjekker jeg straks DR eller en anerkendt avis, før jeg overhovedet overvejer at dele noget.}{eksperter \ipa{ɛɡˈsbɛɐ̯dɐ} · anerkendt \ipa{ˈænɐˌkʰɛnˀd}}{好的，首先我会看是谁发布了这篇文章或帖子，以及是否具名引用了具体专家或官方机构。如果信源匿名，或者语气显得十分愤怒和片面，我在考虑转发前一定会立刻去 DR 或权威报纸核实。}
\dialogturn{3}{Eksaminator}{Kan du give nogle eksempler på, hvordan nyhedsstrømmen på sociale medier påvirker din egen hverdag?}{nyhedsstrømmen \ipa{ˈnyːheðsˌsdʁømˀən} · påvirker \ipa{ˈpʰʌˌviɐ̯ɡə}}{你能举例说明社交媒体上的新闻流如何影响你自己的日常生活吗？}
\dialogturn{4}{Kursist}{Et godt eksempel er, at jeg tidligere tjekkede telefonen lige før sengetid og ofte faldt i søvn med en følelse af uro over alle verdens kriser. Efter at jeg har fjernet nyhedsapps fra forsiden af min skærm, har jeg fået langt bedre ro i hovedet om aftenen.}{sengetid \ipa{ˈsɛŋəˌtˢiðˀ} · uro \ipa{ˈuːˌʁoːˀ}}{一个很好的例子是，过去我常在睡前刷手机，结果经常带着对世界各种危机的焦虑入睡。自从把新闻应用从手机首屏移除后，我晚上的思绪明显平静多了。}
\dialogturn{5}{Eksaminator}{Hvad er efter din mening fordelene og ulemperne ved, at de fleste læser nyheder gratis på mobilen i stedet for i en trykt avis?}{fordelene \ipa{ˈfɒːˌdeːˀlənə} · ulemperne \ipa{ˈuːˌlɛmbɐnə}}{在你看来，大多数人用手机免费看新闻而非读纸质报纸有哪些利弊？}
\dialogturn{6}{Kursist}{Fordelen er høj tilgængelighed: alle har råd til at følge med, uanset økonomi, og man får straks besked ved vigtige begivenheder. Ulempen er derimod, at medierne får færre abonnenter. Som følge heraf bliver der mindre tid til dybdegående journalistik, og mange læser kun overskrifterne.}{tilgængelighed \ipa{ˈtˢilˌɡɛŋˀəliheðˀ} · abonnenter \ipa{ɑboˈnɛndɐ}}{优点是极高的可及性：无论经济条件如何人人都看得起，且重大事件能即时获悉。然而缺点是媒体订户减少；由此导致深度调查报道时间缩水，许多人只看不深究标题。}
\dialogturn{7}{Eksaminator}{Hvad synes du om den rolle, som public service-medier som DR og TV 2 spiller i Danmark? Hvorfor?}{public service \ipa{ˈpʰʌbliɡ ˈsɜːvis} · rolle \ipa{ˈʁʌlə}}{你如何评价 DR 和 TV 2 等公共服务媒体在丹麦扮演的角色？为什么？}
\dialogturn{8}{Kursist}{Jeg synes, de spiller en helt central rolle for sammenhængskraften i demokratiet. Fordi DR ikke styres af reklamer alene, har vi stadig et fælles udgangspunkt med faktatjekkede nyheder og debat, som alle borgere kan stole på.}{demokratiet \ipa{demokʰʁɑˈtˢiːˀəð} · fælles \ipa{ˈfɛlˀəs}}{我认为它们在维护民主社会凝聚力方面起着核心作用。因为 DR 不单受商业广告驱动，我们才依然拥有经过事实核查的新闻与公共讨论基准，让全体公民能够信赖。}

\begin{transbox}
\small
\textbf{\color{NCDNavy}[Emneopgave 3B 口语陈述示范：Uddannelse: livslang læring i Danmark｜161 ord · 约 1.5 分钟]}\\
\textit{Jeg har valgt at tale om uddannelse i Danmark med fokus på livslang læring \ipa{ˈliːu̯slɑŋˀ ˈlɛːɐ̯eŋ}. På det danske arbejdsmarked forventer man ikke længere, at man er færdig med at lære, den dag man afslutter sin grunduddannelse. Tværtimod tilbydes der mange gode muligheder for efteruddannelse \ipa{ˈɛfdɐˌuðˌdænˀəlsə}, for eksempel AMU-kurser, VUC og lokale aftenskoler, ofte med økonomisk støtte fra arbejdsgiveren eller fagforeningen. Fordelen ved dette velfungerende system er, at medarbejderne løbende får opdateret deres faglige kompetencer \ipa{kʰʌmbəˈtˢɑŋsɐ} inden for nye teknologier som kunstig intelligens, hvilket både styrker den enkeltes jobtryghed og samfundets konkurrenceevne. Udfordringen er dog, at det i praksis kan være ret svært at finde tid og overskud til studier ved siden af et krævende fuldtidsjob og travlt familieliv. For mit eget vedkommende er danskuddannelsen \ipa{ˈdænsɡˌuðˌdænˀəlsən} et konkret bevis på livslang læring i voksenalderen, hvor jeg både lærer det danske sprog og forstår samfundet bedre. Fremadrettet håber jeg fortsat at tage nye kurser, fordi nysgerrighed og nye færdigheder gør hele arbejdslivet langt mere spændende og meningsfuldt.}\\[4pt]
\textbf{\color{NCDTeal}[参考译文]}：我选择谈论丹麦的教育，重点是终身学习。在丹麦劳动力市场上，人们不再指望完成基础学业那天就停止学习。相反，社会提供大量优质的继续教育机会，例如 AMU 技能培训、VUC 成人教育与本地夜校，通常还有雇主或工会提供经济资助。这套成熟体制的优点在于员工能持续更新在人工智能等新技术领域的专业能力，这既增强个人职业安全感，也提升国家竞争力。不过挑战在于，在要求颇高的全职工作与忙碌的家庭生活之外，实际很难挤出时间和精力学习。就我个人而言，成人丹麦语教育（DU3）就是成年阶段终身学习的切实体现，在这里我既学习丹麦语又更深入地理解社会。展望未来，我希望继续修读新的课程，因为保持好奇心与学习新技能会让整个职业生涯精彩且有意义得多。
\end{transbox}

\paragraph*{Emneopgave 3B Del 2 考官追问与高分应答示范（4 题）：}
\dialogturn{1}{Eksaminator}{Vil du fortælle lidt mere om, hvorfor livslang læring er blevet så vigtig på det danske arbejdsmarked?}{arbejdsmarked \ipa{ˈɑːbɑi̯dsˌmɑːɡəð} · vigtig \ipa{ˈveɡdi}}{你能再多谈谈为什么终身学习在丹麦职场变得如此重要吗？}
\dialogturn{2}{Kursist}{Ja, teknologien og særligt kunstig intelligens udvikler sig så hurtigt i disse år, at opgaver i både kontorfag og håndværk hele tiden ændrer sig. Desuden skal danskerne arbejde i flere år end tidligere, og derfor er man nødt til at forny sin viden undervejs for at følge med.}{udvikler \ipa{ˈuðˌveɡlɐ} · viden \ipa{ˈviːðən}}{是的，近年来科技尤其是人工智能发展极快，无论是办公室岗位还是技工行业的工作内容都在不断变化。此外丹麦人的法定工作年限比过去更长，因此必须中途不断更新知识才能跟上时代。}
\dialogturn{3}{Eksaminator}{Kan du give nogle eksempler på, hvad du selv har lært som voksen i Danmark, og hvordan du har brugt det?}{eksempler \ipa{eɡˈsɛmˀblɐ} · voksen \ipa{ˈvʌɡsən}}{你能举几个你在丹麦成年后学到的东西以及如何运用的例子吗？}
\dialogturn{4}{Kursist}{Ud over det danske sprog har jeg lært meget om den digitale hverdag i Danmark, for eksempel at bruge MitID, Digital Post og nye digitale arbejdsredskaber. Det har gjort mig meget mere selvstændig i kontakten med både lægen, kommunen og mine kollegaer.}{selvstændig \ipa{ˈsɛlˌsdɛnˀdi} · redskaber \ipa{ˈʁɛðˌsɡæːˀbɐ}}{除了丹麦语之外，我还学到了很多关于丹麦数字化日常的知识，比如使用 MitID、Digital Post 以及新型数字办公工具。这让我在与医生、市政和同事打交道时独立得多。}
\dialogturn{5}{Eksaminator}{Hvad er fordelene og ulemperne ved at tage efteruddannelse online sammenlignet med fysisk undervisning?}{sammenlignet \ipa{ˈsɑmənˌliːˀnəð} · undervisning \ipa{ˈɔnɐˌviːˀsneŋ}}{与线下实体教学相比，在线参加继续教育有哪些优缺点？}
\dialogturn{6}{Kursist}{Fordelen ved onlinekurser er fleksibiliteten, fordi man slipper for transport og kan læse, når børnene er lagt i seng. Ulempen er imidlertid, at man mangler den direkte samtale og sparring med andre kursister i pauserne, hvor man ofte lærer næsten lige så meget.}{fleksibiliteten \ipa{flɛɡsibiliˈtˢeːˀdən} · sparring \ipa{ˈsbɑːeŋ}}{网课的优点在于灵活，因为省去通勤且能在孩子睡下后再学习。然而缺点在于缺少了课间与其他学员的直接交流和思想碰撞——而课间交流往往同样让人受益匪浅。}
\dialogturn{7}{Eksaminator}{Hvad synes du om mulighederne for voksen- og efteruddannelse i Danmark sammenlignet med i dit hjemland? Hvorfor?}{mulighederne \ipa{ˈmuːliˌheːˀðɐnə} · hjemland \ipa{ˈjɛmˌlanˀ}}{与你母国相比，你如何评价丹麦的成人与继续教育机会？为什么？}
\dialogturn{8}{Kursist}{Jeg synes, mulighederne i Danmark er helt unikke, fordi staten støtter så bredt økonomisk, og fordi mange arbejdsgivere ser efteruddannelse som en gevinst frem for en udgift. I mit hjemland må voksne ofte betale dyre kurser helt selv ved siden af lange arbejdsdage.}{unikke \ipa{uˈniɡə} · gevinst \ipa{ɡeˈvensd}}{我觉得丹麦的机会非常独特，因为国家在经济上支持面极广，而且许多雇主将员工进修视为收益而非单纯开支。在我母国，成年人往往只能在超长工时之外自费高价报班。}

\clearpage
\paragraph*{Mundtlig Opgave 2 Del 1（个人看图必答三题示范 · Kursist A \& Kursist B）：}

\begin{grammatikbox}{Kursist A 个人应答示范（图片 A：长者求助阅读 Digital Post｜含 IPA 与译文）}
\small
\begin{enumerate}
  \item \textbf{Hvorfor tror du, det er sådan?} \textit{Jeg tror, mange ældre borgere har brug for hjælp, fordi de offentlige systemer ændrer sig hurtigt, og ikke alle er vokset op med computere og smartphones. Nogle er også usikre på sikkerheden \ipa{ˈseɡɐheːðˀən} eller bange for at trykke forkert, når de skal åbne vigtige breve fra kommunen eller sygehuset.}（我认为许多老年市民需要帮助，是因为公共数字系统变化很快，并非人人都伴随电脑和智能手机长大。有些人也对网络安全心存顾虑，或在打开市政厅和医院重要信件时害怕点错。）
  \item \textbf{Hvilke fordele kan der være ved, at kontakten med det offentlige foregår digitalt?} \textit{Den store fordel er, at sagerne behandles \ipa{beˈhænˀləs} meget hurtigere, fordi breve kommer frem på få sekunder i stedet for at vente på posten. For de fleste borgere er det også praktisk at kunne ordne skat, flytning eller lægetider hjemmefra hele døgnet uden at stå i kø.}（最大优势在于事务处理速度快得多，信件几秒钟内就能送达而不必等邮差。对大多数市民而言，全天候在家就能办理税务、迁址或预约医生而不必排队也十分方便。）
  \item \textbf{Mener du, der kan være nogen ulemper ved, at næsten alt foregår digitalt? Hvorfor? / Hvorfor ikke?} \textit{Ja, den største ulempe er, at sårbare grupper som ældre eller folk med sprogvanskeligheder risikerer at blive hægtet af. Hvis man mister sit MitID eller får problemer med telefonen, kan man pludselig føle sig helt magtesløs \ipa{ˈmɑɡdəsˌløːˀs}, og derfor skal der altid findes personlig hjælp.}（是的，最大缺点是长者或存在语言障碍的脆弱群体可能掉队。一旦丢失 MitID 或手机故障，人会突然感到束手无策，因此必须始终保留人工协助渠道。）
\end{enumerate}
\end{grammatikbox}

\begin{grammatikbox}{Kursist B 个人应答示范（图片 B：咖啡馆里的低头族朋友｜含 IPA 与译文）}
\small
\begin{enumerate}
  \item \textbf{Hvorfor tror du, de gør det?} \textit{Jeg tror, mange automatisk tjekker mobilen af ren vane \ipa{ˈvæːnə}, fordi apps er designet til hele tiden at sende notifikationer og fange vores opmærksomhed. Nogle bruger måske også skærmen til hurtigt at slå noget op sammen eller svare på en praktisk besked hjemmefra.}（我认为许多人纯粹出于习惯下意识看手机，因为应用的设计初衷就是不停发推送来抢占注意力。也有些人可能只是借屏幕快速共同查询信息或回复家里的事务消息。）
  \item \textbf{Hvilke fordele kan der være ved altid at være online på mobilen?} \textit{Fordelen er, at man altid kan få fat i sine nærmeste, finde vej på kortet og tjekke åbningstider eller rejseplaner på få sekunder. Samtidig er det nemt at dele billeder og koordinere \ipa{kʰoʌdiˈneːˀɐ} aftaler med venner i sidste øjeblik, hvilket giver stor frihed i en travl hverdag.}（优点是随时能联系上亲友，几秒内查看地图导航、营业时间或出行路线；同时分享照片和临时与朋友协调聚会也极方便，在忙碌日常中带来很大自由。）
  \item \textbf{Mener du, der kan være nogen ulemper ved, at mobilen fylder så meget i sociale situationer? Hvorfor? / Hvorfor ikke?} \textit{Ja, helt klart! Når alle sidder i deres egen digitale verden på en café, mister man det ægte nærvær \ipa{ˈnɛːɐ̯ˌvɛːˀɐ} og den dybe samtale ansigt til ansigt. Vennerne kan føle, at man ikke rigtig lytter eller prioriterer dem, hvilket på længere sigt kan skade venskabet.}（是的，非常明显！当大家在咖啡馆里各自沉浸在数字世界时，人就会丧失真正的陪伴感与面对面的深入交谈。朋友们会觉得你没有真正倾听或重视彼此相聚，长远来看会伤害友谊。）
\end{enumerate}
\end{grammatikbox}

\clearpage
\paragraph*{Mundtlig Opgave 2 Del 2（Paropgave：丹麦公共政务全数字化 · 8 轮双人辩论示范，兼顾个人与社会维度）：}

\dialogturn{1}{Kursist A}{Det her er jo noget, vi alle oplever i hverdagen i Danmark! Personligt synes jeg, at MitID og Digital Post har gjort rigtig meget nemmere for den enkelte borger. Man kan betale regninger, se prøvesvar fra lægen og ændre adresse hjemme fra sofaen på fem minutter. Hvad synes du?}{hverdagen \ipa{ˈvɛːɐˌdæːˀən} · nemmere \ipa{ˈnɛmˀɐə}}{这确实是我们大家在丹麦日常都会经历的事！我个人觉得 MitID 和 Digital Post 为普通市民带来了极大便利，坐在沙发上五分钟就能缴账单、看医生化验单或改住址。你觉得呢？}
\dialogturn{2}{Kursist B}{Jeg er enig i, at det sparer utrolig meget tid for den enkelte, når man først kender systemet. Samtidig føles det ret sikkert, at alle officielle breve ligger samlet ét sted og ikke bare forsvinder i en postkasse. Men tror du ikke, der er nogle, som har svært ved at følge med?}{officielle \ipa{ʌfiˈɕɛlˀə} · postkasse \ipa{ˈpʰʌsdˌkʰɑsə}}{我同意一旦熟悉系统后，这确实能为个人省下海量时间。而且所有官方信件集中保存、不会在实体信箱丢失也让人觉得挺安全。不过你不觉得还是有些人很难跟上吗？}
\dialogturn{3}{Kursist A}{Jo, desværre. Særligt ældre mennesker, borgere med ordblindhed eller nyankomne udlændinge kan hurtigt føle sig stressede over de mange digitale koder og skemaer. Som følge heraf bliver nogle afhængige af hjælp fra deres børn eller naboer. Men hvad tænker du om fordelene for hele samfundet?}{afhængige \ipa{ˈɑu̯hɛŋˀiə} · samfundet \ipa{ˈsɑmˌfɔnˀəð}}{是的，很遗憾确实如此。尤其是老年人、读写障碍者或新来丹麦的外国人，很容易对五花八门的数字验证码和表格感到焦虑。因此部分人不得不依赖子女或邻居帮忙。不过你怎么看它对全社会层面的好处？}
\dialogturn{4}{Kursist B}{For samfundet giver digitaliseringen jo enorme besparelser på papir, porto og administration, hvorved flere offentlige ressourcer i stedet kan bruges på sygehuse, skoler og ældrepleje, hvilket også gavner klimaet. Men hvad ser du som de største ulemper for samfundet som helhed?}{digitaliseringen \ipa{diɡitˢɑliˈseːˀɐeŋən} · besparelser \ipa{beˈsbɑːˀəlsɐ}}{对整个社会而言，数字化在纸张、邮费和行政人工上节省了巨额开支，从而让更多公共资源得以投入医院、学校和养老护理，而且对气候也很友好。但你觉得从全社会角度看最大的弊端是什么？}
\dialogturn{5}{Kursist A}{Den største risiko for samfundet er nok sårbarheden over for nedbrud og cyberangreb. Hvis MitID eller netværket går ned i blot et halvt døgn, stopper både lægebesøg, betalinger og offentlige ydelser. Derudover skaber det en ny form for social ulighed, hvis man straffes med gebyrer eller forsinkelser, bare fordi man ikke har den nyeste telefon.}{sårbarheden \ipa{ˈsɒːbɑːˌheːðˀən} · cyberangreb \ipa{ˈsɑi̯bɐˌænˌɡʁeːˀb}}{对社会而言最大的风险大概在于面对系统瘫痪和网络攻击时的脆弱性。如果 MitID 或网络哪怕宕机半天，看病、支付和公共福利发放都会停摆。此外，如果一个人只因没有最新款手机就面临滞纳或延误，这也会制造新的社会不平等。}
\dialogturn{6}{Kursist B}{Det er en rigtig vigtig pointe. Jeg hørte angiveligt også i radioen, at flere ældre ikke tør åbne nye links af frygt for svindel, så digital tillid kan ikke tages for givet. Synes du så, vi skal gå tilbage til papirbreve til alle, eller hvordan kan det løses bedre?}{angiveligt \ipa{ænˈɡiːˀvəlið} · svindel \ipa{ˈsvenˀəl}}{这一点切中要害。我似乎也在广播里听说，不少老人因为怕遇到诈骗连新链接都不敢点开，可见数字信任绝不能想当然。那你觉得我们应该全员退回纸质信件吗，还是有更好的解决办法？}
\dialogturn{7}{Kursist A}{Nej, vi skal bestemt ikke droppe de digitale løsninger, for de fungerer fremragende for flertallet. Men kommunerne bør gøre det endnu lettere at få gratis hjælp på bibliotekerne og sikre, at man altid kan ringe til et rigtigt menneske uden at vente en time i telefonkø. Hvad mener du om den balance?}{fremragende \ipa{ˈfʁɑmˌʁɑːˀənə} · telefonkø \ipa{tˢeləˈfoːˀnˌkʰøːˀ}}{不，我们当然不该放弃数字方案，因为它对绝大多数人来说运作得非常出色。但市政部门应当让人更容易在图书馆获得免费协助，并确保随时都能打通真人客服而不必在电话队列苦等一小时。你怎么看这种平衡？}
\dialogturn{8}{Kursist B}{Jeg er helt enig! Hvis vi kombinerer smarte digitale løsninger med god personlig service og nem adgang til fritagelse for dem, der har brug for det, får både den enkelte borger og hele samfundet det bedste fra begge verdener.}{fritagelse \ipa{ˈfʁiˌtˢæːˀəlsə} · borger \ipa{ˈbɒːˀɐ}}{我完全同意！只要把聪明的数字方案、优质的人工服务以及对有需要者的便捷豁免结合起来，市民个人和整个社会就能兼得两全其美之效。}
